\section{Introduction}
    \label{sec:intro}
Echocardiography is a vital imaging modality for cardiac assessment, valued for its non-invasive nature and real-time imaging capabilities. Accurate segmentation of cardiac chambers from echocardiography videos is a foundational step for quantitative analysis~\cite{el2022semi}. Delineating the left ventricle is crucial for calculating the left ventricular ejection fraction, a cornerstone metric for diagnosing cardiovascular diseases and guiding treatment strategies \cite{amer2021resdunet,yang2024echocardiography}.
%
                \input{fig/challenge/cha}
%

However, achieving precise segmentation in echocardiography faces severe challenges stemming from poor image quality and complex cardiac dynamics. \Cref{fig:cha:1,fig:cha:2} show that ultrasound images are characterized by high speckle noise and low contrast, which obscure tissue structures and lead to weak or incomplete anatomical boundaries. These artifacts hinder the model's ability to learn robust features and can result in the inaccurate segmentation. In the temporal dimension, the heart undergoes significant non-rigid deformation throughout the cardiac cycle. \Cref{fig:cha:3,fig:cha:4,fig:cha:5,fig:cha:6} show that, the shape and scale of the LV change dramatically between systole and diastole. This substantial dynamic variation demands that segmentation models possess strong temporal modeling capabilities to accurately track the changing appearance of the target across the video sequence~\cite{li2024automatic}.

Convolutional Neural Networks~\cite{oshea2015introductionconvolutionalneuralnetworks}, particularly U-Net architectures \cite{ronneberger2015u,long2015fully,iglovikov2018ternausnet}, marked a significant milestone in echocardiography segmentation. Their proficiency in learning hierarchical spatial features improved delineation robustness against noise~\cite{jha2019resunet++}. Yet, their inherently local receptive fields limit performance where boundaries are ambiguous and preclude the direct modeling of inter-frame temporal dependencies~\cite{painchaud2022echocardiography}. Vision Transformers address the local-context limitation by providing a global field of view \cite{wang2025transformer,liao2023left}, but their quadratic computational complexity and large data requirements present practical barriers. Spatiotemporal models incorporating recurrent units like ConvLSTM~\cite{chen2021assessing} aim for temporal coherence, yet they risk propagating errors through sequences and add significant computational overhead. A gap thus remains for a model that can integrate global temporal context efficiently while preserving local spatial precision.

This paper introduces \textbf{\mymodel}, a new architecture for echocardiography video segmentation designed to synthesize global temporal modeling with local feature precision in a computationally efficient manner. The core of \mymodel is a Linear Key-Value Association (LKVA) mechanism that captures inter-frame relationships. 
To efficiently update and preserve essential temporal information, we propose Gated Delta Rule (GDR) module that discards irrelevant historical memory while retaining features pertinent to the current frame. 
Key-Pixel Feature Fusion (KPFF) module enhances spatiotemporal representation by integrating local key features, global key features, and pixel-level features.
We evaluate \mymodel on two widely adopted echocardiography video datasets, CAMUS~\cite{leclerc2019deep} and EchoNet-Dynamic~\cite{ouyang2019echonet}, and compare it with various task-specific state-of-the-art methods as well as recent STM models. Experimental results indicate that \mymodel achieves consistently higher segmentation accuracy than existing approaches, demonstrating its effectiveness and robustness on both datasets.
Our main contributions are summarized as follows:
\begin{itemize}
  \item We propose a novel model \mymodel for echocardiography video segmentation, which fully leverages the powerful representational capacity of linear key-value association with new adaptations to ultrasound videos.
  \item We introduce Gated Delta Rule (GDR) and Key-Pixel Feature Fusion (KPFF) module; GDR employs interactions between the current frame and memory bias, helping the model manage its memory. KPFF enables the model to acquire global spatiotemporal features. 
  \item We conduct extensive experiments on the CAMUS and EchoNet-Dynamic datasets to validate the superior performance of our approach over current state-of-the-art methods.
\end{itemize}